\chapter{Експерименти и оценка}
\label{ch:eksperimenti}

Тази глава представя експерименталната постановка, използваните метрики и
резултатите. Изложението следва нарастващата сложност на моделите от
глава~\ref{ch:metod}: таблична базова линия, детерминистичен трансформър и
такива с вероятностни глави, последвани от проверка за точност
при нечуван изпълнител (умение за генерализация на домейна) и
анализ на интерпретируемостта и грешките.

\section{Експериментална постановка}
\label{sec:postanovka}

Всички модели се обучават върху обучаващото подмножество на E-GMD, преценява се
коя да е последната епоха на трениране чрез оценка по валидационното подмножество
и се оценяват финално върху тестовото подмножество, като
се използва общ измервателен модул — така числата на различните модели са пряко
сравними. Освен официалното разделяне (което споделя всеки барабанист между
подмножествата) се провежда и вторична проверка с премахване от обучението на файловете,
изсвирени от конкретен барабанист (раздел~\ref{sec:holdout}).
Възпроизводимостта е фиксирана със seed~$42$.

Табличният модел (LightGBM) е обучен с ранно спиране (най-добра итерация:~$206$-та).
Трансформърът има входен размер на енкодера $d_{\text{model}}=128$, $4$ слоя, $8$ глави, прозорец с
дължина $512$. Обучение посредством оптимизатор \eng{AdamW}~\cite{kingma2015adam}
с $\text{lr}=3\cdot10^{-4}$ и $L1$ функция на загубата.;
обучен е на Apple MPS устройство с достигната най-добра епоха~$14$-та.
Вероятностните глави са инициализирани от запазения обучен трансформър и дообучени с NLL
(функция на загубата, отмерваща отрицателната логаритмична вероятност).

\section{Метрики за оценка}
\label{sec:metriki}

Средната абсолютна грешка сама по себе си е недостатъчна: модел може да я
минимизира като осреднява изводите и така заличава варияциите в динамиката,
което е обратно на целта на задачата ни. Затова се използва набор от
допълващи се метрики:
\begin{itemize}
  \item \textbf{MAE} (средна абсолютна грешка) и \textbf{RMSE} (средноквадратична грешка) — абсолютна точност.
  \item \textbf{Корелация по запис} (Pearson $r$ и Spearman $\rho$) — вярна ли е
        \emph{формата} на динамичната крива.
  \item \textbf{Корелация в рамките на такта} ($\rho$) — правилно ли са подредени
        акцентите и призрачните ноти.
  \item \textbf{Отношение на стандартните отклонения} (\eng{std ratio}) — запазва
        ли моделът естествения обхват на динамиката или го свива (пряка мярка за
        „осредняването''). Стойност $1$ е идеална.
  \item \textbf{Разстояние на Васерщайн} ($W_1$) и \textbf{сечение на хистограмите}
        (\eng{histI}, като $1$ = идентични разпределения) — колко близко семплираният
        изход възпроизвежда истинското разпределение на \eng{velocity}.
  \item \textbf{Дискретна NLL} (\eng{disc\_nll}) — отрицателна логаритмична
        правдоподобност над дискретизирана числова величина, сравнима между всички глави.
\end{itemize}
Като базови летви за сравнение служат глобалната средна динамика (\eng{velocity}) и
най-вече \emph{справочна таблица} (\eng{lookup table}) от средни стойности по
(глас $\times$ жанр $\times$ метрична позиция) — „наивният предсказател на динамика''.

\section{Резултати}
\label{sec:rezultati}

\subsection{Табличен модел}
\label{sec:rez-a}

Табл.~\ref{tab:plan_a} обобщава резултатите на табличния модел спрямо базите за сравнение.
LightGBM \emph{побеждава и двете базови линии по всяка ос}: по-ниско MAE
($-3{,}4$ спрямо справочната таблица) и RMSE, и значително по-висока корелация
както по запис ($0{,}71$ срещу $0{,}55$), така и в рамките на такта ($0{,}67$
срещу $0{,}56$). Глобалната средна не създава никаква вариация (std ratio $0$),
докато LightGBM възстановява по-голямата част от динамичния обхват.

\begin{table}[H]
  \centering
  \caption{Тестово подмножество: табличният модел срещу ,,наивните'' базови линии.}
  \label{tab:plan_a}
  \begin{tabular}{@{}lrrrrrr@{}}
    \toprule
    модел & MAE & RMSE & $r$ (запис) & $\rho$ (запис) & $\rho$ (такт) & std ratio \\
    \midrule
    глобална средна   & 29,67 & 34,64 & — & — & — & 0,000 \\
    справочна таблица  & 21,39 & 28,38 & 0,552 & 0,512 & 0,561 & 0,729 \\
    \textbf{LightGBM} & \textbf{18,02} & \textbf{24,07} & \textbf{0,706} & \textbf{0,665} & \textbf{0,672} & 0,692 \\
    \bottomrule
  \end{tabular}
\end{table}

Анализът на важността на характеристиките (фиг.~\ref{fig:lgbm_importance}) показва, че
най-силният единичен предиктор е стилът (жанр/подстил), следван от темпото
и от времето от последния удар на същия глас. Тоест това
колко отдавна не е свирил \emph{този} ударен инструмент предсказва силата му по-добре,
за разлика от това колко отдавна не е свирил \emph{който и да е} ударен инструмент —
позволяващо моделиране на акцентни цикли на фуса, призрачни ноти на соло барабана, и т.н.

\begin{figure}[H]
  \centering
  \includegraphics[width=0.85\textwidth]{fig_lgbm_importance.png}
  \caption{Важност на признаците (LightGBM), топ~20. Доминират стилът,
  темпото и времевите разстояния в рамките на гласа, следвани от метричната позиция и
  локалната плътност.}
  \label{fig:lgbm_importance}
\end{figure}

По жанрове трудността варира почти двукратно (MAE от $\approx11{,}8$ до
$\approx24{,}1$). Поп музиката е най-трудният жанр (най-непредсказуема
динамика), докато жанрове със силни идиоматични акцентни конвенции
(госпъл, пънк, Ню Орлиънс, реге) са най-лесни.

\subsection{Невронна мрежа с трансформър архитектура}
\label{sec:rez-b}

Табл.~\ref{tab:plan_b} добавя трансформъра. Той \emph{побеждава} базовите летви по всяка ос.
Спрямо LightGBM има малко по-висока грешка по MAE и RMSE, по-ниска корелация, но пък за сметка
на това запазва значително по-голяма част от динамичния обхват на записите.

\begin{table}[H]
  \centering
  \caption{Тестово подмножество: трансформър спрямо предишните модели.
  $\Delta_{\text{std}}$ = средна абсолютна разлика между предсказаното и истинското
  стандартно отклонение на изпълнение (по-ниско е по-добре).}
  \label{tab:plan_b}
  \begin{tabular}{@{}lrrrrrrr@{}}
    \toprule
    модел & MAE & RMSE & $r$ & $\rho$ & $\rho$ (такт) & std ratio & $\Delta_{\text{std}}$ \\
    \midrule
    глобална средна   & 29,67 & 34,64 & — & — & — & 0,000 & — \\
    справочна таблица  & 21,39 & 28,38 & 0,552 & 0,512 & 0,561 & 0,729 & 8,28 \\
    \textbf{LightGBM} & \textbf{18,02} & \textbf{24,07} & \textbf{0,706} & \textbf{0,665} & \textbf{0,672} & 0,692 & 9,02 \\
    трансформър        & 19,56 & 26,82 & 0,605 & 0,556 & 0,612 & \textbf{0,895} & \textbf{5,47} \\
    \bottomrule
  \end{tabular}
\end{table}

Оста, по която трансформърът ясно печели, е точно онази, която средната грешка в точка
скрива — \emph{запазването на обхвата}. Неговите изходи задържат $\approx90\%$
от истинския динамичен обхват (std ratio $0{,}895$ срещу $0{,}692$ на LightGBM), а
грешката в стандартното отклонение по изпълнения е почти наполовина ($5{,}47$ срещу
$9{,}02$). Това потвърждава хипотезата, че модел може да спечели по
MAE ,,изкуствено'' и да изгуби динамиката на изпълненията.
За задача по възстановяване на динамиката, при която вариацията \emph{е} човешкото,
това е истинско предимство и силна мотивация за експериментиране с вероятностните глави.

\begin{figure}[H]
  \centering
  \includegraphics[width=0.75\textwidth]{fig_transformer_training.png}
  \caption{Валидационно MAE по епохи на обучението на трансформъра ($45{,}8 \to 19{,}0$, най-добра
  епоха~$14$).}
  \label{fig:transformer_training}
\end{figure}

\subsection{Вероятностни изходни глави}
\label{sec:rez-c}

Както бе упоменато в раздел~\ref{ch:metod}, за последен слой на мрежата биват изпробвани
и три различни вероятностни глави: гаусова, MDN и категорийна.
Отделно, изходът от всяка глава се третира и оценява по два начина -
\emph{точков} (детерминистичен) и \emph{семплиран} (вероятностен). Точковият изход е
средната стойност на предсказаното разпределение, докато семплираният изход се получава
от въпросното разпределение.

Основният резултат тук е, че семплирането възстановява динамиката: при
всяка от трите изпробвани глави \emph{семплираният} изход има значително по-високо 
отношение на стандартното отклонение (std ratio) от \emph{точковия}, за сметка на 
малко по-голяма средна абсолютна грешка (MAE).
Табл.~\ref{tab:plan_c} дава пълното сравнение.

\begin{table}[H]
  \centering
  \caption{Тестово подмножество: вероятностни глави при точков и семплиран изход,
  спрямо предходните модели. Функцията \eng{disc\_nll} е сравнима между всички глави
  (по-ниско е по-добре).}
  \label{tab:plan_c}
  \begin{tabular}{@{}lrrrrrrr@{}}
    \toprule
    модел / изход & disc\_nll & MAE & $r$ & $\rho$ (такт) & std ratio & $W_1$ & histI \\
    \midrule
    справочна таблица       & — & 21,39 & 0,552 & 0,561 & 0,729 & — & — \\
    LightGBM       & — & \textbf{18,02} & \textbf{0,706} & \textbf{0,672} & 0,692 & — & — \\
    трансф. (точков)   & — & 19,56 & 0,605 & 0,612 & 0,895 & — & — \\
    \midrule
    трансф. + гаусова гл. (точков)         & 3,55 & 22,02 & 0,575 & 0,616 & 0,564 & — & — \\
    трансф. + гаусова гл. (семпл.)          & 3,55 & 27,62 & 0,293 & 0,356 & 0,851 & 6,03 & 0,873 \\
    трансф. + MDN гл. (точков)             & 3,13 & 19,69 & 0,617 & 0,608 & 0,807 & — & — \\
    трансф. + MDN гл. (семпл.)              & 3,13 & 24,55 & 0,455 & 0,468 & \textbf{0,995} & 2,95 & 0,921 \\
    трансф. + катег. гл. (точков) & \textbf{3,04} & 19,17 & 0,632 & 0,626 & 0,784 & — & — \\
    трансф. + катег. гл. (семпл.)  & \textbf{3,04} & 24,16 & 0,468 & 0,483 & 0,982 & \textbf{1,49} & \textbf{0,947} \\
    \bottomrule
  \end{tabular}
\end{table}

От резултатите се открояват няколко извода:
\begin{enumerate}
  \item LightGBM остава ненадминат по средна грешка и корелация, но със свит
    обхват на динамиката.
  \item Семплирането възстановява обхвата, но за сметка на точковата точност.
  \item Категорийната глава печели по повечето други съществени метрики:
    най-добро NLL ($3{,}04$), най-добро точково MAE сред невронните модели, и
    най-съответстващо на разпределението при семплиране (std ratio $0{,}982$,
    $W_1=1{,}49$, histI $0{,}947$). Тъй като \eng{velocity} е ограничено цяло
    число, категорийната дискретна глава не ,,пилее'' вероятностна маса извън
    $[0,127]$. Семплираният изход на категорийната глава съвпада с истинското
    разпределение почти точно (фиг.~\ref{fig:categorical_hist}).
  \item MDN надминава гаусовата и се доближава повече до категорийната с
    доста подобни метрики.
\end{enumerate}

\begin{figure}[H]
  \centering
  \includegraphics[width=0.7\textwidth]{fig_categorical_hist.png}
  \caption{Семплирано срещу истинско разпределение (в тестовото подмножество) на силата
  (\eng{velocity}) за категорийната глава).
  Наслагването е видимо най-плътното сред всички глави (histI $0{,}947$).}
  \label{fig:categorical_hist}
\end{figure}

\subsection{Устойчивост към нечуван изпълнител}
\label{sec:holdout}

Всички досегашни числа използват официално публикуваното разделяне на набора данни,
в което всеки барабанист примсъства измежду всички подмножества.
За да проверим умението за генерализация на обучените модели към \emph{нечуван} изпълнител,
преразделяме данните така, че двама барабанисти (drummer3 и drummer8) със значими на брой
записи и стилове да образуват тестово множество, което никога не се появява по време на обучението,
при запазване на всички $17$ жанра в обучението. Гръбнакът (трансформърът) и категорийната глава се
обучават наново от нулата върху това разделяне.

Резултатът е реална разлика в генерализацията (табл.~\ref{tab:plan_d}):
точковото MAE скача от $\approx19$ на $\approx28$ ($+45\%$), а семплираното
разстояние на Васерщайн експлодира от $1{,}49$ на $16{,}82$. Това \emph{не е}
артефакт от по-малкото данни: при преобучението валидационната MAE върху
\emph{видени} барабанисти достигна $19{,}05$, което съвпада с представянето в
разпределението.

\begin{table}[H]
  \centering
  \caption{Устойчивост към задържан барабанист (drummer3 + drummer8) спрямо
  представянето в разпределението.}
  \label{tab:plan_d}
  \begin{tabular}{@{}llrrrrr@{}}
    \toprule
    модел / изход & подмножество & MAE & $r$ & $\rho$ (такт) & std ratio & $W_1$ \\
    \midrule
    гръбнак (точков) & старо тестово\    & 19,56 & 0,605 & 0,612 & 0,895 & — \\
    гръбнак (точков) & ново тестово & 27,65 & 0,478 & 0,491 & 0,837 & — \\
    категор.\ (точков) & старо тестово\   & 19,17 & 0,632 & 0,626 & 0,784 & — \\
    категор.\ (точков) & ново тестово & 27,98 & 0,490 & 0,489 & 0,706 & — \\
    категор.\ (семпл) & старо тестово\    & 24,16 & 0,468 & — & 0,982 & 1,49 \\
    категор.\ (семпл) & ново тестово & 31,96 & 0,309 & — & 0,921 & 16,82 \\
    \bottomrule
  \end{tabular}
\end{table}

В основата на проблема стои това, че абсолютната сила на удара е
индивидуална черта на изпълнителя. Задържаните извън обучаващото подмножество
барабанисти просто удрят по-силно: средна сила (\eng{MIDI velocity}) $84{,}6$ срещу $63{,}2$ за обучаващия набор, а
моделът предсказва $\approx67{,}7$ — по същество наученото от обучението, изместено
ниско (фиг.~\ref{fig:velocity_shift}). Изместването от $-16{,}8$ единици почти
точно съвпада със семплирания $W_1$ от $16{,}82$: „правилна форма, грешно място''.
Разлагайки точковата разлика, около $3{,}8$ MAE е чисто абсолютно отместване
(коригираната за отместване MAE е $24{,}20$), а останалите $\approx5$ MAE са
истинска загуба на структура в рамките на записа.

\begin{figure}[H]
  \centering
  \includegraphics[width=0.8\textwidth]{fig_velocity_shift.png}
  \caption{Истински разпределения на \eng{velocity}: обучаващ набор срещу задържани
  барабанисти, с предсказаната от модела средна. Изместването \emph{е} разликата в
  генерализацията.}
  \label{fig:velocity_shift}
\end{figure}

Важното е \emph{какво все пак генерализира}: относителната динамика частично се
пренася — корелацията по запис ($0{,}49$) и в рамките на такта ($0{,}49$) остават
ясно положителни, а семплирането все още възстановява обхвата (std ratio $0{,}92$).
Моделът продължава да подрежда акценти срещу призрачни ноти за нов изпълнител; той
просто калибрира погрешно \emph{абсолютното} ниво. Можем да заключим, че
моделът пренася относителната изразителна форма към невидени изпълнители,
но не и абсолютното им динамично ниво. (Тук трябва да отбележим, че
това ново тестово подмножество е нарочно трудно — двамата избрани барабанисти са
сред най-силно удрящите в E-GMD.)

\section{Интерпретируемост и анализ на грешките}
\label{sec:interpretiruemost}

Накрая изследваме дали моделът учи \emph{музикална} структура, или само напасва
числа, и къде греши. Анализът е върху категорийната глава по тестовото подмножество.

\paragraph{Ембединги на каноничен глас и жанр}
Наученото пространство на гласовете ($8$-мерно) откроява различните инструментални
семейства: всички разновидности на удари по фуса се групират, а проекцията разделя
силните гласове (акцентиран соло барабан, чинели) от
тихите (соло барабан, перваз, каса) — моделът е научил \emph{динамичната
идентичност на инструмента} (фиг.~\ref{fig:voice_embedding}).

\begin{figure}[H]
  \centering
  \includegraphics[width=0.8\textwidth]{fig_voice_embedding.png}
  \caption{2D проекция на наученото ембединг пространство на каноничните гласове, оцветено по
  средна сила (\eng{velocity}).}
  \label{fig:voice_embedding}
\end{figure}

Жанровият ембединг
($16$-мерен) е по-слаб и кодира по-скоро сходно \emph{динамично поведение},
отколкото жанрова таксономия (напр.\ сходство на кънтри със соул).

\begin{figure}[H]
  \centering
  \includegraphics[width=0.8\textwidth]{fig_genre_embedding.png}
  \caption{2D проекция на наученото ембединг пространство на жанровете, оцветено по
  средна сила (\eng{velocity}).}
  \label{fig:genre_embedding}
\end{figure}

\paragraph{Акценти и призрачни ноти}
Моделът възпроизвежда почти точно метричната йерархия: нотите на силно време (\eng{on-beat})
са по-силни от тези на слабо време (\eng{off-beat}), фиг.~\ref{fig:metrical}
(силно време: истинско $75{,}9$ / предсказано $75{,}0$; слабо време: $58{,}4$ / $57{,}4$).

\begin{figure}[H]
  \centering
  \includegraphics[width=0.8\textwidth]{fig_metrical.png}
  \caption{Средна \eng{velocity} по метрична позиция, истинска срещу предсказана.}
  \label{fig:metrical}
\end{figure}

Обикновеният соло барабан е силно многомодален: $50\%$ от истинските удари са под \eng{velocity}~$40$
(призрачни ноти). Точковият изход не улавя това — само $27\%$ от прадсказаните така ноти
са със сила под $40$. Семплирането възстановява локалния максимум на призрачните ноти до
$41$ и запазва до голяма степен формата на разпределението (видно на фиг.~\ref{fig:ghost_notes}). 
Точковата оценка повишава значително силата на призрачните ноти до такава степен, че са
на прага да бъдат въобще наричани ,,призрачни''. Семплирането ги връща обратно в реалистичния диапазон.

\begin{figure}[H]
  \centering
  \includegraphics[width=0.75\textwidth]{fig_ghost_notes.png}
  \caption{Хистограма на истинската и предсказана сила (точков извод) на соло барабана.
  С пунктир е указан локалният максимум на призрачните ноти ($40$).}
  \label{fig:ghost_notes}
\end{figure}

Групирайки стойностите за сила в тестовото подмножество
(\eng{true velocity}), виждаме, че отклонението на разпределението на точковия изход
е монотонно и голямо в краищата:
$+32{,}6$ за най-тихите (сила $0$--$11$) и $-27{,}8$ за най-силните (сила $116$--$127$), с
преминаване през нулата около $55$ (фиг.~\ref{fig:dynamic_regression}). Точковият
модел издърпва тихите нагоре, а силните надолу към условното средно — точно
затова ,,свива'' динамиката.

\begin{figure}[H]
  \centering
  \includegraphics[width=0.8\textwidth]{fig_dynamic_regression.png}
  \caption{Отклонение на точковия изход спрямо очакваната сила (групирана по
  дискретни категории).}
  \label{fig:dynamic_regression}
\end{figure}

Грешката следва \emph{колко променлив} е гласът, а не
колко често се среща: най-лесни са гласове с постоянна сила (акцентиран соло
барабан $11{,}1$, каса $15{,}1$), най-трудни — изразителните гласове, служещи предимно
за маркиране на ритъма (камбана на райд $28{,}4$, перваз $26{,}4$, фусове). По метрична позиция
грешката е практически равна (на силно време $19{,}13$, на слабо време $19{,}19$) —
моделът е равномерно компетентен във времето, а уменията му варират по
\emph{инструмент}, не по позиция в такта.

\paragraph{Систематични отмествания по глас и връзка с употребата на различно програмирани бленди
от използвания електронен комплект барабани.}
Открояват се две систематични отклонения по гласове: затвореният фус се предсказва
твърде тихо ($-13{,}8$), а первазът — твърде силно ($+14{,}5$). Тези
отмествания до голяма степен са \emph{артефакт на данните}. Рендирането през 43
комплекта пренасочва падове към различни гласове според пресета
(раздел~\ref{sec:egmd}), така че характеристиката „глас'' носи шум и моделът учи
\emph{смесеното} средно. Затова погласова рекалибрация би просто ,,излекувала симптом'',
но не и същината на проблема, което би било подвеждащо и често невалидно за други бленди и комплекти
барабани. Вместо това,
правилната корекция е на ниво данни (преизграждане върху една бленда и/или стандартизация) и е
очертана сред насоките за бъдещо развитие (раздел~\ref{sec:budeshti}).
